\newcommand*{\tW}{\widetilde{W}}

Пусть~$A$ --- частично упорядоченное множество.

Определим функцию~$S$ её действием на множествах: $S(A) := A \cup \set*{A}$. Назовём множество~$S(A)$ множеством, \emph{следующим за~$A$}. Тогда

\begin{align*}
	0 &:= \emptyset,
	\\
	1 &:=
	S(0) =
	S(\emptyset) =
	\emptyset \cup \set*{\emptyset} =
	\set*{\emptyset},
	\\
	2 &:= S(1) =
	S \bigl( S(\emptyset) \bigr) =
	\set*{\emptyset} \cup \set*{\set*{\emptyset}} =
	\set*{\emptyset, \set*{\emptyset}},
	\\
	3 &:= S(2) =
	S \Bigl( S \bigl( S(\emptyset) \bigr) \Bigr) =
	\set*{\emptyset, \set*{\emptyset}}
	\cup
	\set*{\set*{\emptyset, \set*{\emptyset}}}
	=
	\set*{
		\emptyset, \set*{\emptyset}
		\set*{\emptyset, \set*{\emptyset}}
	}
	,
\end{align*}
и так далее.

Сложение чисел определяется тогда как объединение множеств:
\[
	1 + 2
	:=
	\set{1, 2} =
	\set*{\emptyset} \cup
	\set*{\emptyset, \set*{\emptyset}}
	=
\]

У нас есть конструкция натуральных чисел, начиная \emph{с нуля} (ноль мы считаем натуральным числом, следуя Н.\,Бурбаки). Теперь определим множество целых чисел~$\ZZ$.

Интуитивно хочется получить что-то, состоящее из нуля~$\set*{0}$, положительных натуральных чисел~$\NN \setminus \set*{0}$ и <<отрицательных натуральных чисел>> $-(\NN\setminus\set*{0}) = \set{-n}{n \in \NN \setminus \{0\}}$. Проблема в том, что операция <<минус>> не определена. Можно её определить аксиоматически, но мы пойдём другим путём.

Определим целые числа с помощью пар натуральных чисел. Геометрическая иллюстрация этого процесса показана на~\picref{pic:integers}: мы отождествляем точки вдоль прямых в~$\NN^2$ и <<нанизываем их>> на горизонтальный и вертикальный лучи --- экземпляры~$\NN$. Развернув <<угол>>, получившийся после отождествления точек, получим прямую, на которой живут целые числа.

Две точки $(a, b)$ и $(m, n)$ из~$\NN^2$ лежат на одной из прямых, изображённых на~\picref{pic:integers}, тогда и только тогда, когда $a - m = b - n$ (вектор с координатами $(a - m, b - n)$ коллинеарен направляющему вектору прямой).


Проблема только в том, что для натуральных чисел не определена операция вычитания.

Эта трудность не преодолевается, но обходится: заметим, что равенство $a - m = b - n$ равносильно равенству $a + n = b + m$. Сложение определено для натуральных чисел, так что его и примем за определение того, что $(a, b)$ и $(m, n)$ лежат на одной прямой. Поэтому
$\ZZ = \NN^2 / \sim$, где
\[
	\bigl(
		(a, b) \sim
		(m, n)
	\bigr)
	:=
	\bigl(
		a + n = b + m
	\bigr)
	.
\]
Обозначим класс эквивалентности числа~$(m, n) \in \NN^2$ символом $[(m, n)]$:
\[
	[(m, n)] :=
	\set{(a, b) \in \NN^2}{(a, b) \sim (m, n)}.
\]

При этом положительные числа в~$\ZZ$ эквивалентны точкам на горизонтальной прямой, а отрицательные --- точкам на вертикальной прямой.

% =============================================================================

Заметим, что при $a \in \NN$ и $k \in \NN$
\begin{align*}
	[(a, 0)] &= [(a + k, k)];
	\\
	[(0, a)] &= [(k, a + k)].
\end{align*}
Поэтому если в паре~$(a, b)$ имеем $a > b$, то класс $[(a, b)]$ представляет положительное число, а если $a < b$ --- то отрицательное.


Определим отображение $+ \colon \ZZ \times \ZZ \to \ZZ$. Если, интуитивно, при $a \in \NN$ класс $[(a, 0)] \in \ZZ$ определяет неотрицательное число $a \in \ZZ$, а класс $[(0, a)] \in \ZZ$ --- отрицательное число $-a \in \ZZ$, то определяем


Пока мы не ввели сложения и умножения в~$\NN^2 / \sim$. Подумаем, как можно его ввести, используя знания из школы.

\subsection[Вводим операции в $\NN^2/\sim$]{Вводим операции в $\bm{\NN^2/\sim}$}

\emph{Интуитивно} мы можем понять, что такое число $m - n$ при $m \geq n$. Поэтому при $m \geq n$ имеем $(m, n) \sim (m - n, 0)$ --- и пара $(m, n)$ представляет число $m - n \in \NN$. А при $m < n$ имеем $(m, n) \sim (0, n - m)$.

% =============================================================================

В этом разделе мы опишем аналог векторного пространства (или модуля над кольцом) для $\NN^2$. Но сначала надо с ним познакомиться.

\subsection{Определения множеств с операциями}

\begin{definition}
	Множество~$M$ с бинарной операцией $* \colon M \times M \to M$ называется \tdef{магмой}.
	\label{<+label+>}
\end{definition}
От операции~$*$ требуется только \tdef{замкнутость}: применение операции~$*$ к паре $(m, n) \in M \times M$ должно давать элемент множества~$M$.

\begin{definition}
	Если в магме~$(M, *)$ потребовать от операции~$*$ ещё и ассоциативности: $(a
	* b) * c = a * (b * c)$ для всех $a$, $b$, $c$ из~$M$, --- то получим
	определение \tdef{полугруппы}%
	\footnote{%
		Заметка на полях: как мы узнаем позже, $(\ZZ, +)$ --- это \emph{группа}.
		А~вот если считать, что в~$\NN$ нет нуля, то $(\NN, +)$ ---
		\emph{полу}группа, а её подлежащее множество~$\NN$--- это <<половина
		от~$\ZZ$>>. Возможно, поэтому дали такое название?
	}%
	.
\end{definition}

\begin{definition}
	Если в полугруппе $(M, *)$ есть элемент~$e \in M$, такой, что для любого $a \in M$ выполняются равенства $a * e = e * a = a$, то элемент~$e$ называется \tdef{нейтральным}, а $(M, *)$ --- \tdef{моноидом}.
	\label{<+label+>}
\end{definition}

Итак, множество~$\NN$ с операцией~$+ \colon \NN \times \NN \to \NN$ --- это моноид: сложение ассоциативно, а в качестве нейтрального (по сложению) элемента выступает~$0 \in \NN$.

Далее, множество~$\NN$ с операцией $\cdot \colon \NN \times \NN \to \NN$ --- это тоже моноид: умножение ассоциативно, а нейтральный (по умножению) элемент --- это $1 \in \NN$.

\subsubsection{Группа}
\begin{definition}
	Множество~$(G, *)$ c бинарной операцией $* \colon G \times G \to G$ называется \tdef{группой}, если:
	\begin{enumerate}[1)]
		\setcounter{enumi}{-1}
		\item $G$ замкнуто относительно~$*$;
		\item $*$ ассоциативна: $a * (b * c) = (a * b) * c$;
		\item в~$G$ существует нейтральный относительно~$*$ элемент~$e \in G$, такой, что для любого элемента $a \in G$
			\[
				a * e = e * a = a
				;
			\]
		\item у всякого $a \in G$ существует обратный относительно~$*$ элемент~$\hat{a} \in G$, то есть такой, что
			\[
				\hat{a} * a = e = a * \hat{a}
				.
			\]
	\end{enumerate}
	Если сложение удовлетворяет ещё аксиоме
	\begin{enumerate}
		\setcounter{enumi}{3}
		\item для всяких~$a$, $b \in G$ операция~$*$ коммутативна:
			\[
				a * b = b * a,
			\]
	\end{enumerate}
	то группа называется \tdef{абелевой}.
	\label{<+label+>}
\end{definition}

\subsubsection{Поле}

\begin{definition}
	\tdef{Поле} --- это тройка~$(F, +, \cdot)$, состоящая из множества~$F$ и двух
	операций~$+$ и~$\cdot$ на нём, удовлетворяющая условиям:
	\begin{enumerate}[1)]
		\item $(F, +)$ --- абелева группа с нейтральным элементом~$0$;
		\item $(F\setminus\set*{0}, \cdot)$ --- абелева группа с нейтральным элементом~$1$;
		\item $0 \cdot 1 = 0$;
		\item $a \cdot (b + c) = a\cdot b + a \cdot c$ для любых $a$, $b$ и $c \in F$;
	\end{enumerate}

	\label{<+label+>}
\end{definition}<++>

\subsubsection{Кольцо}

\subsubsection{Модуль над кольцом}

\subsection[Аналог модуля над кольцом для $\NN^2$]{Аналог модуля над кольцом для $\bm{\NN^2}$}

Опишем теперь

\subsubsection{Операции в~$\NN^2$}

Будем рассматривать $W$ как <<модуль>> над полукольцом~$\NN$ (в отличие от <<настоящего>> модуля над кольцом, у нас нет операции $v \mapsto -v$). Пусть $\bas{e} = (\bm{e}_1, \bm{e}_2)$ --- базис в~$W$. Будем считать, что $\bm{e}_1$ направлен вправо, а $\bm{e}_2$ --- вверх, как на~\picref{pic:integers}.

Будем отождествлять элемент $W \ni a\bm{e}_1 + b\bm{e}_2$ с набором координат~$(a, b) \in \NN^2$ в базисе~$\bas{e}$. Поэтому $W$ изоморфно $\NN^2$: изоморфизм устанавливается соответствием
\[
	\NN^2 \ni (a, b) \mapsto a\bm{e}_1 + b \bm{e}_2 \in W
	.
\]
Будем обозначать тот факт, что модуль $W$ изоморфен~$\NN^2$, символом $W \cong \NN^2$.

Сначала определим сложение в~$W \cong \NN^2$: пусть 
\[
	+ \colon \NN^2 \times \NN^2 \to \NN^2
	\colon
	\bigl(
		(a, b),
		(m, n)
	\bigr)
	\mapsto
	\bigl(
		a + m, b + n
	\bigr)
	.
\]
Мотивация для этого --- посмотреть на сложение в~$W$:
\[
	(a\bm{e}_1 + b\bm{e}_2) + (m\bm{e}_1 + n\bm{e}_2) =
	(a + m)\bm{e}_1 + (b + n)\bm{e}_2
	.
\]

Если мы хотим, чтобы на~\picref{pic:integers} вектор~$\bm{e}_1$ соответствовал будущему числу~$1 \in \ZZ$, а вектор --- будущему числу~$-1 \in \ZZ$, то достаточно положить их сумму эквивалентной~$0$: $\bm{e}_1 + \bm{e}_2 \sim 0$. Так как $0 = 0\bm{e}_1 + 0\bm{e}_2$, то в координатах это означает, что $(1, 0) + (0, 1) \sim (0, 0)$.

Теперь перейдём к фактормодулю $\tW := W/\sim$, изоморфному $\NN^2 / \sim =: \ZZ$.

Определим сложение в нём унаследованным из сложения в~$\NN^2$:
\[
	+\colon
	\tW
	\times
	\tW
	\to
	\tW
	\colon
	\bigl(
		[(a, b)],
		[(m, n)]
	\bigr)
	\mapsto
	[(a, b) + (m, n)]
\]
Иными словами, $[(a, b)] + [(m, n)] = [(a, b) + (m, n)]$. 

Для фактормодуля~$\tW$ поэтому имеем
\[
	[v] + [w] := [v+w]
\]
для любых векторов~$v$, $w \in W$.

Теперь определим билинейное умножение на~$\tW$, превратив его в <<алгебру>>. Если хотим уметь вычислять $[v] \cdot [w]$, то, разложив эти векторы по базису: $\bm{v} = v^1 \bm{e}_1 + v^2 \bm{e}_2$, $\bm{w} = w^1 \bm{e}_1 + w^2 \bm{e}_2$, --- получим, что
\[
	[v] \cdot [w] =
	[v^1 \bm{e}_1 + v^2 \bm{e}_2]
	\cdot
	[w^1 \bm{e}_1 + w^2 \bm{e}_2]
	=
	\bigl(
		v^1 [\bm{e}_1] +
		v^2 [\bm{e}_2]
	\bigr)
	\cdot
	\bigl(
		w^1 [\bm{e}_1] +
		w^2 [\bm{e}_2]
	\bigr)
	=
	\bigl(
		v^1 w^1 \,
		[\bm{e}_1] \cdot
		[\bm{e}_1]
		+
		v^1 w^2 \,
		[\bm{e}_1] \cdot
		[\bm{e}_2]
		+
		v^2 w^1 \,
		[\bm{e}_2] \cdot
		[\bm{e}_1]
		+
		v^2 w^2 \,
		[\bm{e}_2] \cdot
		[\bm{e}_2]
	\bigr)
	.
\]
Итак, умножение векторов в~$\tW$ определяется умножением классов $[\bm{e}_\nu]$
базисных векторов.  Теперь, в соответствии с тем, что $[\bm{e}_1]$ --- это
вектор в направлении <<положительных>> чисел, а $[\bm{e}_2]$ --- в направлении
<<отрицательных>>, определим умножение классов базисных векторов:
\[
	\begin{array}{r@{{} = {}}l@{\quad}l}
		[\bm{e}_1] \cdot [\bm{e}_1] & [\bm{e}_1]
		& \text{(<<$1 \cdot 1 = 1$>>);}
		\\
		[\bm{e}_1] \cdot [\bm{e}_2] & [\bm{e}_2]
		& \text{(<<$1 \cdot (-1) = -1$>>);}
		\\
		[\bm{e}_2] \cdot [\bm{e}_1] & [\bm{e}_1]
		& \text{(<<$(-1) \cdot 1 = -1$>>);}
		\\
		[\bm{e}_2] \cdot [\bm{e}_2] & [\bm{e}_1]
		& \text{(<<$(-1) \cdot (-1) = 1$>>).}
	\end{array}
\]

Сопоставляя $\NN^2 / \sim \ni [(a, b)] \longleftrightarrow a[\bm{e}_1] + b [\bm{e}_2] \in \tW$, имеем:
\begin{align*}
	[(a, b)] \cdot [(m, n)] &=
	[a \bm{e}_1 + b \bm{e}_2] \cdot [m \bm{e}_1 + n \bm{e}_2]
	\nlinea{=}
	am [\bm{e}_1] \cdot [\bm{e}_1] + 
	an [\bm{e}_1] \cdot [\bm{e}_2] + 
	bm [\bm{e}_2] \cdot [\bm{e}_1] + 
	bn [\bm{e}_2] \cdot [\bm{e}_2]
	\nlinea{=}
	am [\bm{e}_1] + 
	an [\bm{e}_2] + 
	bm [\bm{e}_2] + 
	bn [\bm{e}_1]
	\nlinea{=}
	(am + bn) [\bm{e}_1] +
	(an + bm) [\bm{e}_2]
	.
\end{align*}
Это соответствует тому, что $[(a, b)] \cdot [(m, n)] = [(am + bm, an + bm)]$.
